\chapter{Introduction}
\label{chap:intro}

Proving that a quantum spin chain is gapped in the thermodynamic limit requires a lower bound on
the energy difference between the first excited state and the ground state that is uniform in the
system size. The \emph{spacetime bootstrap}, also called the \emph{two-purity doubling criterion},
reduces this infinite-volume problem on periodic rings to two finite-size inequalities, as proved
in the main note~\cite{HaldaneNote} and the scope note~\cite{ScopeNote}. The Lean~4 libraries
\texttt{CertifiedGaps} and \texttt{SeedKernel} in \texttt{lean/} formalize this criterion and its
applications to concrete spin chains; each statement below links to the Lean declaration that
proves it.

\section*{The criterion}

For a periodic ring Hamiltonian \(H_L = \sum_{j=1}^{L} h_{j,j+1} + \sum_{j=1}^{L} o_j\) of \(L\)
unit cells, write \(Z_L(\beta) = \Tr e^{-\beta H_L}\) and define the \emph{spatial purity} and
\emph{thermal purity}
\[
  \Ps(n,\beta) = \frac{Z_{2n}(\beta)}{Z_n(\beta)^2}, \qquad
  \Pt(n,\beta) = \frac{Z_n(2\beta)}{Z_n(\beta)^2}.
\]
A thermal purity \(\Pt(L, \beta)\) near one means that the ground state carries almost all of the
Boltzmann weight, which forces a unique ground state and bounds the spectral gap
\(\Delta_L = E_1(L) - E_0(L)\) from below whenever the thermal defect \(1 - \Pt(L, \beta)\) decays
exponentially in \(\beta\). Squaring the Gibbs operator in imaginary time improves \(\Pt\) at fixed
\(L\), but advancing \(L\) and \(\beta\) together requires a spatial transfer operator. If the bond
interaction is reflection positive, meaning that \(-h = \sum_\alpha C_\alpha \otimes
\overline{C_\alpha}\)
in a local basis with a real symmetric on-site term \(o\), then \(Z_L(\beta)\) is a moment sum
\(\sum_i \mu_i^L\) of a positive transfer operator at even lengths \(L\), and both purities improve
quadratically when length and inverse temperature are doubled together. Two purity bounds at a
single \emph{seed} scale \((n_0, \beta_0)\) then yield a uniform lower bound on \(\Delta_L\) for
every even \(L \ge n_0\).

\section*{What is formalized}

The formalization is organized into five parts.
\begin{enumerate}
  \item Chapter~\ref{chap:criterion} proves the abstract two-purity doubling criterion for any
    family of finite spectra whose partition functions satisfy three \emph{spatial moment
    inequalities} (Theorem~\ref{thm:gap-of-seeds}).
  \item Chapter~\ref{chap:chain} extends the criterion to unequal seed defects by iterating the
    exact asymmetric doubling recurrence through rational chain certificates checked in the Lean
    kernel (Theorem~\ref{thm:gap-of-chain-certificate}).
  \item Chapter~\ref{chap:rings} proves that every reflection-positive ring satisfies the spatial
    moment inequalities by constructing Trotterized positive semidefinite transfer matrices and
    taking the continuous-time limit via the Lie--Trotter product formula
    (Theorem~\ref{thm:rp-moment-bounds}), which yields the gap criterion for rings
    (Theorem~\ref{thm:rp-gap}).
  \item Chapters~\ref{chap:heisenberg} and~\ref{chap:models} apply the ring criterion to the
    spin-1 Heisenberg ring and seven further spin-chain families, establishing a uniform gap
    theorem for each model conditional only on its two seed inequalities.
  \item Chapter~\ref{chap:seedkernel} proves the seed inequalities inside the Lean kernel for the
    longitudinal-field Ising chain at \(g = 1/4\) and \(h = 1/20\), giving an unconditional gap
    theorem (Theorem~\ref{thm:ising-gap}).
\end{enumerate}

\section*{How to read the status}

A statement marked as formalized has a Lean declaration with that exact content, and a proof marked
as formalized is complete in Lean. The open nodes of the dependency graph are the seed inequalities
of each model, which are finite numerical bounds on rings of a few to a few thousand sites left as
\texttt{sorry} in Lean and certified by computer in \texttt{models/}. All other results depend only
on the standard axioms \texttt{propext}, \texttt{Classical.choice}, and \texttt{Quot.sound}.

\begin{remark}[Conventions]
  \label{rem:conventions}
  The spectral gap is stated through the predicate \(\GapGT(e, c)\) (Definition~\ref{def:gapgt}),
  which asserts that the eigenvalue family \(e\) has a nondegenerate minimum (attained at a single
  index) and every other eigenvalue
  lies more than \(c\) above it. This formulation avoids sorting eigenvalues and records
  ground-state uniqueness and the gap lower bound simultaneously. Partition functions of concrete
  rings are defined as \(Z_L(\beta) = \operatorname{Re}\Tr e^{-\beta H_L}\), the real part of a
  complex matrix trace. Ring lengths count unit cells.
\end{remark}
